\chapter{Bronchitis}
\index{Bronchitis}

\section*{Definition and Overview}
Bronchitis is inflammation of the larger airways of the lungs, the bronchi, and its principal symptom is a cough, often productive of phlegm (sputum). Two distinct conditions share the name. Acute bronchitis is a short-lived illness, usually caused by a viral infection, that begins like a cold and resolves within a few weeks. Chronic bronchitis is a long-term condition defined by a productive cough lasting at least three months of the year for two consecutive years; it is a form of chronic obstructive pulmonary disease (COPD) and is almost always caused by smoking. Although they share a name, the two conditions differ substantially in cause, outlook, and treatment, and distinguishing between them is important.

\section*{Epidemiology}
Acute bronchitis is extremely common, affecting people of all ages and peaking in autumn and winter when respiratory viruses circulate. Most adults experience an episode every few years, and it is one of the most frequent causes of short-term illness and time off work. Chronic bronchitis is strongly linked to smoking and is most common in people over 40 who have smoked for many years; it has traditionally been more common in men, although the gap has narrowed as smoking rates in women rose. Long-term exposure to dust, chemical fumes, and air pollution also increases the risk. Chronic bronchitis is a leading cause of chronic illness, disability, and hospital admission in older adults worldwide.

\section*{Aetiology and Pathophysiology}
\begin{itemize}
    \item \textbf{Viral infection:} the great majority of cases of acute bronchitis are caused by the same viruses responsible for colds and influenza, which spread to the bronchi and inflame them
    \item \textbf{Bacterial infection:} less common; may follow a viral illness, particularly in people with chronic lung disease or impaired immunity
    \item \textbf{Smoking:} the principal cause of chronic bronchitis; tobacco smoke irritates the airway lining and, over many years, damages it permanently
    \item \textbf{Inhaled irritants:} dust, chemical fumes, and air pollution contribute to airway inflammation, especially in smokers
    \item \textbf{Mechanism:} the bronchial lining becomes swollen and produces excess mucus; in acute bronchitis this change is temporary and reversible, whereas in chronic bronchitis the mucus-secreting glands enlarge, the airway walls thicken, and the damage becomes irreversible
\end{itemize}

\section*{Clinical Features}
\begin{itemize}
    \item A cough, the cardinal symptom of both forms
    \item Production of phlegm, which may be clear, white, yellow, or green
    \item In acute bronchitis: a runny or blocked nose, mild sore throat, low-grade fever, and malaise, following the pattern of a cold
    \item Chest discomfort or a raw, burning sensation from coughing
    \item Wheeze or a rattling sound in the chest in some people
    \item Shortness of breath, usually mild in acute bronchitis but persistent and progressive in chronic bronchitis
    \item In chronic bronchitis: a daily cough with phlegm lasting months at a time, typically worse in the morning, with recurrent chest infections and exertional breathlessness
\end{itemize}

\section*{Differential Diagnosis}
\begin{itemize}
    \item \textbf{Common cold:} an upper respiratory tract infection with nasal congestion and sore throat but little or no cough
    \item \textbf{Pneumonia:} infection of the air sacs, with higher fever, tachypnoea, breathlessness, and sometimes pleuritic chest pain
    \item \textbf{Asthma:} intermittent wheeze and cough, often triggered by allergens or exercise, without persistent daily sputum
    \item \textbf{COPD:} fixed airflow obstruction, usually in a smoker, of which chronic bronchitis is one component
    \item \textbf{Pertussis (whooping cough):} paroxysms of coughing that may end in a whoop or vomiting
    \item \textbf{Heart failure:} breathlessness and cough from pulmonary congestion, more likely in older people with cardiac disease
    \item \textbf{Lung cancer:} a new or changing cough, haemoptysis, or weight loss, particularly in smokers
    \item \textbf{Gastro-oesophageal reflux:} acid reflux can produce a chronic cough, often worse at night or after meals
\end{itemize}

\section*{When to Seek Medical Care}
See a doctor if a cough lasts more than three weeks, if blood is coughed up, if chest infections recur or worsen, or if breathlessness or a cough does not settle as expected.

\textbf{Contact your local emergency services for:}
\begin{itemize}
    \item Difficulty breathing, fast breathing, or chest pain
    \item Coughing up blood
    \item A high fever, confusion, or extreme drowsiness
    \item Blue or grey lips, or the feeling of being unable to get enough air
    \item Worsening symptoms in someone with heart disease, COPD, or a weakened immune system
\end{itemize}

\section*{Diagnosis and Investigations}
\begin{itemize}
    \item \textbf{History:} the pattern and duration of the cough, sputum production, smoking history, and whether symptoms followed a cold
    \item \textbf{Examination:} chest auscultation, temperature, respiratory rate, and pulse oximetry
    \item \textbf{Chest X-ray:} arranged when pneumonia, heart failure, or another cause is suspected, especially in smokers or older people
    \item \textbf{Lung function tests:} spirometry is the key investigation for chronic bronchitis and COPD, demonstrating airflow obstruction
    \item \textbf{Blood and sputum tests:} occasionally needed in severe, persistent, or atypical cases
    \item \textbf{Review:} follow-up to confirm the cough has settled, particularly in smokers with prolonged symptoms
\end{itemize}

\section*{Management}
\begin{itemize}
    \item \textbf{Acute bronchitis:} usually requires no specific treatment; rest, adequate fluids, and simple analgesia such as paracetamol for fever and discomfort
    \item \textbf{Cough medicines:} not generally recommended; honey in warm water may soothe older children and adults
    \item \textbf{Antibiotics:} not usually helpful because most cases are viral; considered for people with chronic lung disease or significant risk factors
    \item \textbf{Inhalers:} bronchodilator therapy may help wheeze but is not routine in acute bronchitis
    \item \textbf{Chronic bronchitis:} smoking cessation is the single most important intervention; inhalers, pulmonary rehabilitation, and regular exercise control symptoms and reduce exacerbations
    \item \textbf{Vaccination:} annual influenza and pneumococcal vaccination to lower the risk of secondary infection
    \item \textbf{Managing exacerbations:} early use of prescribed medicines and prompt medical review when symptoms worsen
\end{itemize}

\section*{Complications}
\begin{itemize}
    \item \textbf{Pneumonia:} the most important complication of acute bronchitis, especially in older adults, smokers, and those with chronic illness; acute bronchitis may also trigger an asthma attack
    \item \textbf{Recurrent chest infections:} frequent in chronic bronchitis, with each episode potentially worsening lung function
    \item \textbf{Progressive COPD:} chronic bronchitis can advance to severe disease with low oxygen levels, right-heart strain, and respiratory failure
    \item \textbf{Long-term disability:} breathlessness may limit work, activity, and quality of life
    \item \textbf{Smoking-related harm:} an increased risk of lung cancer and cardiovascular disease in people with chronic bronchitis
\end{itemize}

\section*{Prognosis and Outcomes}
Acute bronchitis almost always resolves within two to three weeks, although a dry cough can persist for several more weeks; it causes no lasting damage in otherwise healthy people. Chronic bronchitis is a lifelong condition in which the airway damage cannot be reversed, but stopping smoking markedly slows the decline in lung function and reduces exacerbations. With good treatment and self-management, many people with chronic bronchitis remain active for many years, although advanced disease can cause substantial breathlessness and disability.

\section*{Prevention and Self-Care}
\begin{itemize}
    \item Do not smoke, and avoid second-hand smoke and lung irritants at home and at work
    \item Wash hands regularly and avoid close contact with people with coughs and colds
    \item Keep influenza and pneumococcal vaccination up to date
    \item Rest, drink plenty of fluids, and use simple analgesia during acute bronchitis
    \item Do not take antibiotics unless prescribed by a doctor
    \item Stay physically active, and consider pulmonary rehabilitation for chronic bronchitis
    \item Manage coexisting conditions such as asthma, diabetes, and heart disease
    \item See a doctor for any cough lasting more than three weeks, or for new or changing symptoms
\end{itemize}

\section*{Sources and Further Reading}
The following organisations provide reliable information on bronchitis for patients and healthcare students.

\begin{itemize}
    \item NHS. \textit{Bronchitis: symptoms, treatment, and when to see a doctor.} https://www.nhs.uk/conditions/bronchitis/
    \item Mayo Clinic. \textit{Bronchitis: causes, diagnosis, and treatment.} https://www.mayoclinic.org/
    \item MedlinePlus. \textit{Bronchitis health information.} https://medlineplus.gov/bronchitis.html
    \item MSD Manual. \textit{Acute and chronic bronchitis: professional overview.} https://www.msdmanuals.com/
    \item Centers for Disease Control and Prevention. \textit{Chest cold (acute bronchitis) information for patients.} https://www.cdc.gov/
    \item National Institute for Health and Care Excellence. \textit{Respiratory infection and COPD guidance.} https://www.nice.org.uk/
\end{itemize}
